\documentclass[11pt,a4paper]{article}

% ---------- encoding & lingua ----------
\usepackage[utf8]{inputenc}
\usepackage[T1]{fontenc}
\usepackage[italian]{babel}
\usepackage{lmodern}

% ---------- layout ----------
\usepackage[a4paper,margin=2.4cm,headheight=14pt]{geometry}
\usepackage{microtype}
\usepackage{setspace}
\onehalfspacing
\usepackage{parskip}

% ---------- grafica ----------
\usepackage{graphicx}
\usepackage{xcolor}
\usepackage{tikz}
\usepackage{booktabs}
\usepackage{array}
\usepackage{caption}
\captionsetup{font=small,labelfont=bf,width=0.95\textwidth}
\usepackage{enumitem}
\setlist{itemsep=2pt,topsep=3pt}
\usepackage{amsmath}
\usepackage{amssymb}

% ---------- intestazioni ----------
\usepackage{fancyhdr}
\pagestyle{fancy}
\fancyhf{}
\fancyhead[L]{\small\itshape Metriche di valutazione --- Trentingrana Captioning}
\fancyhead[R]{\small\thepage}
\renewcommand{\headrulewidth}{0.4pt}

% ---------- titoli ----------
\usepackage{titlesec}
\definecolor{prim}{HTML}{2C5F8A}
\definecolor{gold}{HTML}{E0A458}
\definecolor{green}{HTML}{5B8C5A}
\definecolor{red}{HTML}{C1492E}
\definecolor{boxbg}{HTML}{F2F6FA}
\definecolor{boxbd}{HTML}{C7D6E5}
\titleformat{\section}{\Large\bfseries\color{prim}}{\thesection}{0.6em}{}
\titleformat{\subsection}{\large\bfseries\color{prim!85!black}}{\thesubsection}{0.6em}{}
\titlespacing*{\section}{0pt}{14pt}{6pt}

% ---------- link ----------
\usepackage[colorlinks=true,linkcolor=prim,citecolor=prim,urlcolor=prim]{hyperref}

% ---------- box riassuntivo ----------
\newcommand{\keybox}[2]{%
\begin{center}
\begin{tikzpicture}
\node[draw=boxbd,fill=boxbg,rounded corners=4pt,inner sep=8pt,
      text width=0.93\textwidth,align=left] {%
\textbf{\color{prim}#1}\\[2pt]\small #2};
\end{tikzpicture}
\end{center}}

% scorciatoia figure metriche
\newcommand{\mfig}[1]{figures/metrics/#1}

% ============================================================
\begin{document}

% ---------------- TITOLO ----------------
\begin{center}
{\color{prim}\rule{\textwidth}{1.5pt}}\\[0.5cm]
{\LARGE\bfseries Metriche di valutazione per il captioning sensoriale}\\[0.3cm]
{\large Oltre BLEU: CIDEr, BERTScore, conformità lessicale e CLIPScore}\\[0.4cm]
{\color{prim}\rule{\textwidth}{1.5pt}}\\[0.6cm]
{\small Progetto \textbf{AI4FQC} --- Project 07: \textit{GRANA Captioning} \quad$\bullet$\quad Relazione tecnica complementare}
\end{center}

\vspace{0.3cm}
\keybox{In sintesi}{%
Affianchiamo a BLEU sei metriche complementari. Tre risultati emergono dai dati:
(1) \textbf{BLEU inganna} --- una baseline che ripete sempre la stessa frase
ottiene il BLEU-1 più alto di tutti; (2) \textbf{CIDEr discrimina} i modelli là
dove BLEU li appiattisce, e ordina correttamente \mbox{m1\,$<$\,m3\,$<$\,m6};
(3) \textbf{CLIPScore} --- l'unica metrica che guarda l'immagine --- mostra che
i modelli addestrati \emph{non} battono la baseline costante, confermando in modo
indipendente che, con encoder congelato, l'immagine non viene davvero sfruttata.}

% ============================================================
\section{Perché non basta BLEU}

Nel captioning esistono molti modi corretti di descrivere la stessa immagine, ma
il dataset fornisce \textbf{un solo riferimento} per immagine (la frase del
panelista). Una metrica che confronta parola per parola penalizza le
riformulazioni corrette:

\begin{center}
\small
\begin{tabular}{@{}p{3.2cm} p{10.5cm}@{}}
\toprule
Riferimento & \textit{``Il formaggio ha pasta compatta con occhiatura rada.''}\\
Predizione & \textit{``Il formaggio ha pasta dura con piccola occhiatura.''}\\
\midrule
\multicolumn{2}{@{}p{13.7cm}@{}}{\small Stesso significato, parole diverse.
BLEU-4 $\approx 0{,}1$ (nessuna sequenza di 4 parole condivisa), ma la
descrizione è \textbf{corretta}.}\\
\bottomrule
\end{tabular}
\end{center}

Per questo usiamo sette metriche, divise in tre famiglie:
\begin{itemize}
\item \textbf{Basate sul testo} (predizione vs riferimento testuale):
BLEU-1/4, METEOR, ROUGE-L, CIDEr, BERTScore.
\item \textbf{Basata sull'immagine} (predizione vs immagine, ignora il
riferimento): CLIPScore.
\item \textbf{Specifica del dominio}: Conformità Vocabolario.
\end{itemize}

Tutti i numeri provengono da \texttt{eval\_metrics/metrics\_summary.csv}: 43 file
di predizioni (4 baseline + 3 modelli), su 7 attributi per-attributo e 3 versioni
\emph{global}.

% ============================================================
\section{Le sette metriche, una per una}

\subsection{BLEU-1 e BLEU-4}
\textbf{Cosa misura.} La precisione degli $n$-gram: quante sequenze di $n$ parole
consecutive della predizione compaiono nel riferimento. BLEU-1 (singole parole)
$\to$ ``ho usato le parole giuste?''; BLEU-4 (4 parole) $\to$ ``ho usato la
fraseologia giusta, nell'ordine giusto?''.

\textbf{Range e lettura.} Valore in $[0,1]$, più alto è meglio. Per questo
dominio: BLEU-4 $>0{,}4$ ottimo; $0{,}2$--$0{,}4$ discreto; $<0{,}15$ scarso.

\textbf{Risultati.} Modelli per-attributo BLEU-4 $\approx 0{,}30$--$0{,}45$;
modelli \emph{global} $\approx 0{,}13$ (crollo netto). \emph{Il tranello:} la
baseline \texttt{most\_frequent} emette \emph{sempre} la stessa caption e ottiene
il BLEU-1 più alto di tutti ($0{,}84$--$0{,}90$). BLEU premia la ripetizione del
template, non la comprensione.

\subsection{METEOR}
\textbf{Cosa misura.} Come BLEU-1 ma più indulgente: allinea le parole con
\textbf{stemming} (\textit{compatto}\,$\approx$\,\textit{compatta}) e
\textbf{sinonimi}, bilanciando precisione e recall.

\textbf{Range e lettura.} $[0,1]$; tende a essere più alto di BLEU-4.
$>0{,}6$ forte; $0{,}4$--$0{,}6$ parziale; $<0{,}35$ divergente.

\textbf{Risultati.} Per-attributo $\approx 0{,}53$--$0{,}66$; global $\approx
0{,}29$. Racconta la stessa storia di BLEU $\to$ buon segnale di robustezza.

\subsection{ROUGE-L}
\textbf{Cosa misura.} La più lunga sottosequenza comune (LCS): parole presenti in
entrambe le frasi nello stesso ordine, non necessariamente consecutive. Premia il
rispetto dell'ordine generale senza richiedere contiguità.

\textbf{Range e lettura.} $[0,1]$; tipicamente tra BLEU-4 e METEOR.

\textbf{Risultati.} Per-attributo $\approx 0{,}50$--$0{,}67$; global $\approx
0{,}30$. Si muove in tandem con METEOR/BLEU.

\subsection{CIDEr}
\textbf{Cosa misura.} È la metrica pensata \emph{apposta} per il captioning. Usa
$n$-gram pesati con \textbf{TF-IDF}: i termini rari e informativi
(\textit{occhiatura}, \textit{piccante}) valgono molto, quelli comuni
(\textit{il}, \textit{ha}) quasi nulla. Risponde a: ``il modello azzecca i
termini sensoriali \textbf{salienti}, o riempie con parole generiche?''.

\textbf{Range e lettura.} $\ge 0$, senza tetto fisso (scala $\times 10$). Va letto
\textbf{in relazione tra modelli sullo stesso attributo}. Implementazione
standalone a \emph{un} riferimento: confrontabile fra i nostri modelli, non con
benchmark esterni a riferimenti multipli (es.\ COCO).

\textbf{Risultati.} CIDEr aggiunge informazione vera: \emph{discrimina} dove BLEU
appiattisce. La media per-attributo ordina i modelli correttamente
($\text{random } 0{,}23 < \text{freq\_w } 0{,}37 < \text{m1 } 0{,}59 <
\text{m3 } 0{,}65 < \text{m6 } 0{,}70$): i modelli addestrati battono nettamente
le baseline casuali, e \mbox{m6} è il migliore. \texttt{most\_frequent} resta in
testa ($0{,}79$) solo perché il riferimento \emph{è} spesso la frase frequente
(Fig.~\ref{fig:cider}).

\subsection{BERTScore}
\textbf{Cosa misura.} La similarità \textbf{semantica} via embedding contestuali
(BERT italiano): non le parole esatte, ma i significati. Riportiamo l'F1.
Risponde a: ``il modello ha capito il \textbf{concetto}, anche riformulando?''.

\textbf{Range e lettura.} Per l'italiano tipicamente $[0{,}7,\,1{,}0]$: lo ``zero''
pratico è alto. Va letto in modo relativo: $>0{,}90$ forte; $0{,}84$--$0{,}90$
buono; $<0{,}83$ distante.

\textbf{Risultati.} Tutti i valori in $0{,}83$--$0{,}92$: poco discriminante,
perché tutte le frasi parlano di formaggio con la stessa struttura. Utile come
\emph{sanity check}, non per il ranking. I modelli global scendono a $\approx
0{,}80$ (ancora i più deboli).

\subsection{Conformità Vocabolario (metrica custom)}
\textbf{Cosa misura.} La percentuale di parole della caption che appartengono al
\textbf{vocabolario controllato} del progetto (1278 forme sensoriali attestate).
Risponde a una domanda che nessuna metrica standard pone: ``il modello «parla
formaggio» col registro certificato, o inventa termini?''.

\textbf{Range e lettura.} $[0,1]$. Il vocabolario contiene solo lemmi sensoriali
(non parole funzionali \textit{il/ha/di}), quindi anche una caption perfetta non
raggiunge $1{,}0$: leggere in relazione. $>0{,}55$ forte; $0{,}45$--$0{,}55$
moderato; $<0{,}45$ fuori registro.

\textbf{Risultati.} Modelli addestrati $\approx 0{,}48$--$0{,}61$. Più alta sugli
attributi descrittivi (Texture $\approx 0{,}60$--$0{,}67$, Struttura $\approx
0{,}57$). I modelli global hanno conformità pari o superiore ai per-attributo, pur
avendo BLEU/METEOR peggiori: frasi più ``lessicalmente corrette'' ma meno fedeli
al riferimento specifico.

\subsection{CLIPScore (la più importante per il captioning)}
\textbf{Cosa misura.} L'\textbf{unica} metrica che guarda l'immagine. Proietta
immagine e caption nello stesso spazio CLIP (versione multilingue per l'italiano)
e ne calcola la similarità coseno:
\[
\text{CLIPScore} = \cos\!\big(\text{CLIP}_{\text{img}}(\text{fetta}),\;
\text{CLIP}_{\text{txt}}(\text{caption})\big).
\]
\textbf{Ignora il riferimento del panelista.} Risponde alla domanda più onesta:
``questa caption descrive bene \emph{questa immagine}?''. Aggira il problema del
riferimento singolo.

\textbf{Range e lettura.} Per CLIP multilingue la coseno grezza è bassa in
assoluto ($\approx 0{,}15$--$0{,}30$ anche per coppie corrette): \textbf{non è una
percentuale}. Conta il confronto relativo e la direzione. Sul nostro setup:
$\approx 0{,}20$--$0{,}21$ attributo ben descrivibile; $\approx 0{,}18$ poco
legato all'aspetto visivo.

\textbf{Risultati.} Due esiti chiave, discussi nella sezione seguente.

% ============================================================
\section{Risultati trasversali}

\subsection{CLIPScore: i modelli addestrati non battono la baseline costante}

\begin{figure}[ht]
\centering
\includegraphics[width=0.82\textwidth]{\mfig{m_clip_models}}
\caption{CLIPScore medio sui 7 attributi per-attributo. Gli scarti tra modelli
addestrati e baseline costante sono dentro il rumore ($<0{,}004$).}
\label{fig:clipmodels}
\end{figure}

La media CLIPScore è quasi identica per tutti (Tab.~\ref{tab:clipmodels}): i
modelli addestrati (m1/m3/m6) \emph{non} superano \texttt{most\_frequent}, che
spara sempre la stessa frase ignorando l'immagine. \textbf{Interpretazione:} con
l'encoder congelato i modelli non àncorano la caption ai contenuti visivi --- fanno
\emph{language modeling} sulla distribuzione delle caption. È la conferma
quantitativa e indipendente del risultato dello \emph{shuffle test}: se i modelli
usassero l'immagine, batterebbero la baseline costante. Non lo fanno.

\begin{table}[ht]
\centering
\small
\caption{Metriche medie per modello (7 attributi per-attributo). CIDEr ordina i
modelli; CLIPScore resta piatto.}
\label{tab:clipmodels}
\begin{tabular}{@{}l c c c c@{}}
\toprule
Modello & BLEU-4 & CIDEr & Vocab & \textbf{CLIPScore}\\
\midrule
most\_frequent (costante) & 0,362 & 0,790 & 0,524 & \textbf{0,1902}\\
freq\_weighted            & 0,362 & 0,370 & 0,529 & \textbf{0,1919}\\
random                    & 0,358 & 0,234 & 0,527 & \textbf{0,1913}\\
m1\_cnn\_lstm             & 0,384 & 0,586 & 0,535 & \textbf{0,1925}\\
m3\_vit\_transformer      & 0,392 & 0,646 & 0,532 & \textbf{0,1921}\\
m6\_vit\_gpt              & 0,398 & 0,704 & 0,538 & \textbf{0,1933}\\
\bottomrule
\end{tabular}
\end{table}

\subsection{CLIPScore distingue attributi visibili da non visibili}

\begin{figure}[ht]
\centering
\includegraphics[width=0.82\textwidth]{\mfig{m_clip_attr}}
\caption{CLIPScore medio per attributo. Gli attributi visivi (texture, colore,
struttura) ottengono valori più alti; olfatto/gusto restano bassi per chiunque.}
\label{fig:clipattr}
\end{figure}

La variazione vera è \textbf{per attributo}, non per modello. CLIP ``vede'' che
texture ($0{,}208$), colore ($0{,}206$) e struttura ($0{,}197$) sono descrivibili
dall'immagine, mentre aroma ($0{,}181$), sapore ($0{,}184$) e profumo ($0{,}185$)
restano bassi: nessuna immagine può rivelare un sapore. È un controllo di sanità
che \emph{valida} il comportamento della metrica. Eccezione istruttiva:
\textbf{Spessore della Crosta} è fisicamente visibile ma basso ($0{,}184$), perché
la crosta è una porzione piccola della fetta e CLIP la pesa poco.

\subsection{CIDEr discrimina dove BLEU appiattisce}

\begin{figure}[ht]
\centering
\includegraphics[width=0.92\textwidth]{\mfig{m_cider_attr}}
\caption{CIDEr per attributo. A differenza di BLEU-4 (valori vicini fra i
modelli), CIDEr separa i modelli e premia i termini sensoriali rari azzeccati.}
\label{fig:cider}
\end{figure}

Su \emph{Spessore della Crosta}, BLEU-4 dà $0{,}39$--$0{,}44$ per tutti, mentre
CIDEr va da $1{,}04$ (freq\_weighted) a $2{,}29$ (most\_frequent), con i modelli
addestrati in mezzo ($1{,}55$--$1{,}76$): un ordinamento che BLEU non vedeva.

\subsection{Il tranello di BLEU, visualizzato}

\begin{figure}[ht]
\centering
\includegraphics[width=0.95\textwidth]{\mfig{m_bleu_trap}}
\caption{Stessa coppia di modelli (\texttt{most\_frequent} vs m6) vista da due
metriche. A sinistra BLEU-1: la caption costante \emph{vince}. A destra
CLIPScore: i due sono praticamente uguali. BLEU-1 da solo è fuorviante.}
\label{fig:trap}
\end{figure}

\texttt{most\_frequent} ha il BLEU-1 più alto (caption costante che combacia sul
prefisso) ma CLIPScore identico a tutti. Una metrica da sola può mentire: serve
leggerle insieme.

\subsection{Profilo completo di un modello}

\begin{figure}[ht]
\centering
\includegraphics[width=0.92\textwidth]{\mfig{m_heatmap}}
\caption{Profilo di m6 (ViT+GePpeTto) per attributo: valori grezzi annotati,
colore normalizzato per colonna (rende comparabili scale diverse). Si nota la
correlazione tra CIDEr alto e gli attributi a caption dominante (Spessore), e la
relativa uniformità di CLIP.}
\label{fig:heatmap}
\end{figure}

% ============================================================
\section{Come riprodurre}

Ambiente isolato in \texttt{eval\_metrics/.venv} (non tocca le dipendenze del
training).
\begin{enumerate}
\item \texttt{.\textbackslash eval\_metrics\textbackslash setup.ps1}
 --- crea il venv e installa le dipendenze (una tantum).
\item \texttt{python eval\_metrics\textbackslash add\_image\_paths.py}
 --- aggiunge \texttt{image\_path} ai \texttt{predictions.csv} (per CLIPScore).
\item \texttt{python eval\_metrics\textbackslash compute\_metrics.py predictions
 -{}-images-root .} --- calcola tutte le metriche; scrive
 \texttt{metrics\_summary.csv}.
\end{enumerate}

\begin{table}[ht]
\centering
\small
\caption{Implementazione delle metriche.}
\begin{tabular}{@{}l l l@{}}
\toprule
Metrica & Libreria / metodo & Tipo\\
\midrule
BLEU-1/4, METEOR, ROUGE-L & HuggingFace \texttt{evaluate} & testuale, 1 rif.\\
CIDEr & implementazione standalone (TF-IDF $n$-gram) & testuale, 1 rif.\\
BERTScore-F & \texttt{bert-score}, \texttt{lang="it"} & semantica\\
Conformità Vocabolario & custom su \texttt{vocabulary.csv} & dominio\\
CLIPScore & \texttt{open\_clip}, \texttt{xlm-roberta-base-ViT-B-32} (LAION-5B) & image-grounded\\
\bottomrule
\end{tabular}
\end{table}

% ============================================================
\section{Limiti e caveat}
\begin{itemize}
\item \textbf{Riferimento singolo}: tutte le metriche testuali confrontano contro
una sola caption; CIDEr e BLEU ne soffrono strutturalmente. CLIPScore (che ignora
il riferimento) è il complemento più prezioso.
\item \textbf{CIDEr} qui è a un riferimento: confrontabile fra i nostri modelli,
non con benchmark esterni.
\item \textbf{BERTScore} è poco discriminante su questo dominio (frasi simili).
\item \textbf{CLIPScore} usa CLIP generico (non fine-tuned sul formaggio): coglie
tratti grossolani (colore, texture), non sfumature sensoriali fini; i valori
assoluti bassi sono normali.
\item \textbf{CLIPScore costante tra modelli} è esso stesso il risultato: il collo
di bottiglia è l'\textbf{encoder congelato}. Il prossimo esperimento naturale è il
\textbf{fine-tuning dell'encoder} e la ri-misura di CLIPScore: se i modelli si
staccano dalla baseline, l'immagine viene finalmente sfruttata.
\end{itemize}

% ============================================================
\appendix
\section{Tabella completa per attributo}

\begin{table}[ht]
\centering
\small
\setlength{\tabcolsep}{4pt}
\caption{BLEU-4 / CIDEr / CLIPScore per attributo, modelli addestrati vs
\texttt{most\_frequent}. Dati da \texttt{metrics\_summary.csv}.}
\begin{tabular}{@{}l ccc ccc ccc@{}}
\toprule
& \multicolumn{3}{c}{BLEU-4} & \multicolumn{3}{c}{CIDEr} & \multicolumn{3}{c}{CLIPScore}\\
\cmidrule(lr){2-4}\cmidrule(lr){5-7}\cmidrule(lr){8-10}
Attributo & m3 & m6 & m.f. & m3 & m6 & m.f. & m3 & m6 & m.f.\\
\midrule
Aroma     & 0,45 & 0,45 & 0,44 & 0,53 & 0,62 & 0,53 & 0,182 & 0,183 & 0,178\\
Profumo   & 0,38 & 0,38 & 0,32 & 0,35 & 0,39 & 0,33 & 0,183 & 0,186 & 0,187\\
Sapore    & 0,41 & 0,42 & 0,40 & 0,53 & 0,61 & 0,86 & 0,184 & 0,184 & 0,181\\
Texture   & 0,34 & 0,35 & 0,29 & 0,34 & 0,42 & 0,42 & 0,208 & 0,210 & 0,204\\
Spessore  & 0,40 & 0,42 & 0,44 & 1,55 & 1,76 & 2,29 & 0,184 & 0,185 & 0,184\\
Struttura & 0,31 & 0,33 & 0,21 & 0,41 & 0,41 & 0,52 & 0,198 & 0,199 & 0,193\\
Colore    & 0,45 & 0,44 & 0,42 & 0,81 & 0,73 & 0,59 & 0,205 & 0,207 & 0,205\\
\bottomrule
\end{tabular}
\end{table}

\noindent\textit{m.f.} = \texttt{most\_frequent}. La tabella completa con tutte le
8 metriche per tutti i 43 file è in \texttt{eval\_metrics/metrics\_summary.csv}.

\end{document}
